\section{Complexity Analysis}

This appendix analyzes the computational cost of preserving agency.

\subsection{Constraint Detection Cost}

Let $\ConstraintFn : \States \to \{0,1\}$ be nontrivial. Then constraint evaluation is at least as hard as distinguishing semantically adjacent states.

\begin{lemma}
If $\States$ has nonzero semantic dimension, constraint detection is not constant-time.
\end{lemma}

\begin{proof}
Constraint satisfaction depends on global properties of $x \in \States$. Any constant-time procedure would imply collapse of semantic distinctions, contradicting nontriviality.
\end{proof}

\subsection{Compression–Agency Tradeoff}

Let $R$ denote compression ratio and $\Preserve{\AbstractFn}$ constraint preservation.

\begin{theorem}[Compression–Agency Tradeoff]
There exists a constant $K$ such that:
\[
R \cdot \Preserve{\AbstractFn} \le K.
\]
\end{theorem}

\begin{proof}
By the data processing inequality, compression reduces mutual information. Constraint preservation requires mutual information retention. Their product is bounded.
\end{proof}

\subsection{Undecidability of Refusal Detection}

\begin{theorem}
Determining whether refusal is required for arbitrary inputs is undecidable in the general case.
\end{theorem}

\begin{proof}[Sketch]
Constraint evaluation can encode arbitrary halting problems via state encoding. Therefore no total decision procedure exists.
\end{proof}

